\section{The multi-provider datum-difference model}
\label{sec:model}

We model $N$ providers, each of which realizes a lunar reference frame and a lunar
ephemeris. For a single physical body-point $\bm{p}$---a surface asset, or the
lunocentre-referenced Moon state itself---provider $a$ reports a coordinate
$\xstate_a(\bm{p},t)$. The object of study is the pairwise disagreement between two
providers' reports of the \emph{same} physical point.

\subsection{The seven-parameter datum difference}
\label{sec:datumdiff}

To first order, the difference between two providers' coordinates for one point
decomposes into a constant seven-parameter similarity (Helmert) datum and a
time-varying residual,
\begin{equation}
  \xstate_b(\bm{p},t) - \xstate_a(\bm{p},t)
  \;=\; \Jac\!\bigl(\xstate_a\bigr)\,\datum_{ab} \;+\; \resid_{ab}(\bm{p},t),
  \label{eq:diffmodel}
\end{equation}
where $\datum_{ab}\in\R^{7}$ is the constant Helmert datum of $b$ with respect to
$a$, and $\resid_{ab}$ is the residual that no single constant Helmert transform can
absorb---the genuine differential \emph{dynamics}, a prediction disagreement rather
than a coordinate convention. The point-Jacobian is evaluated at the reporting
provider's coordinate $\xstate_a$; since $\xstate_a=\bm{p}+O(\|\datum_{ab}\|)$, we may
substitute the shared physical point $\bm{p}$, $\Jac(\xstate_a)=\Jac(\bm{p})+O(\|\datum_{ab}\|^{2})$,
so at the first order at which \eqref{eq:diffmodel} holds the two are interchangeable.
The decomposition is \emph{unique} exactly when the sampled point geometry gives the
$3\times7$ design a full column rank of seven; for a single lever arm it need not (a
uniform radial scale and a radial translation are then indistinguishable), which is
why the externally validated fit below drops the scale column
(Section~\ref{sec:rotonly}). The Helmert datum collects three translations
$\ttrans=(t_x,t_y,t_z)$, one scale $s$, and three small rotations
$\rot=(\theta_x,\theta_y,\theta_z)$,
\begin{equation}
  \datum \;=\; (t_x,\,t_y,\,t_z,\,s,\,\theta_x,\,\theta_y,\,\theta_z)^\top \in \R^7,
  \label{eq:datumvec}
\end{equation}
and acts on a point $\bm{p}\in\R^3$ through the linearized similarity transform
$\bm{p}\mapsto\bm{p}+\ttrans+s\bm{p}+\rot\times\bm{p}$, whose Jacobian with respect
to $\datum$ is the $3\times7$ point-Jacobian
\begin{equation}
  \Jac(\bm{p}) \;=\; \bigl[\;\bm{I}_3\;\big|\;\bm{p}\;\big|\;-[\bm{p}]_\times\;\bigr],
  \qquad
  [\bm{p}]_\times\,\rot = \bm{p}\times\rot .
  \label{eq:pointjac}
\end{equation}
The first three columns are the translation, the fourth (the point's own
coordinates) is the scale, and the last three are the generators of rotation. This
is the same point-Jacobian that governs the single-provider datum-identifiability
problem of the companion paper \citep{baweja2026datum}; here it acts on the
\emph{difference} between two providers rather than on a single frame's null space.

Given a set of sampled points and epochs, the constant datum $\datum_{ab}$ that best
explains the differences \eqref{eq:diffmodel} is the least-squares solution of the
stacked linear system, and $\resid_{ab}$ is the residual left after removing it. We
solve this by centred normal equations with Jacobi (diagonal) preconditioning, for
the reasons given in Appendix~\ref{app:methods}.

\subsection{The rotation-only sub-fit}
\label{sec:rotonly}

The full seven-parameter fit is used only as an internal analytic check
(Section~\ref{sec:validation}). The externally validated decomposition uses a
\emph{rotation-only} sub-fit: the six-parameter rigid-body (Euclidean) model that
retains translation and rotation but drops the scale column,
\begin{equation}
  \xstate_b(\bm{p},t) - \xstate_a(\bm{p},t)
  \;=\; \ttrans_{ab} \;-\;[\bm{p}]_\times\,\rot_{ab}
   \;+\; \resid^{\mathrm{rot}}_{ab}(\bm{p},t).
  \label{eq:rotonly}
\end{equation}
For a fixed body-point sampled through the physical geometry, the rotation
$\rot_{ab}$ is recovered from the centred point correspondences by the classical
singular value decomposition of their cross-covariance
\citep{arun1987leastsquares}; the residual root-mean-square
$\|\resid^{\mathrm{rot}}_{ab}\|$ measures the part of the disagreement that no
constant translation-plus-rotation removes. We report, for each provider pair, the
raw disagreement, the rotation-fit residual, and the recovered rotation magnitude;
these are the \Validated{} quantities of Section~\ref{sec:interephemeris}. Using the
rigid-body sub-fit rather than the full seven-parameter fit is deliberate: it is the
minimal model whose recovered rotation is unambiguous (a scale column can trade
against a radial translation, whereas a pure rotation cannot), so it is the fit we can
reproduce against an independent oracle to a tight tolerance.

\subsection{Reducible, irreducible, and the interop tax}
\label{sec:defs}

\begin{definition}[Interop tax]
\label{def:tax}
For a user that consumes coordinates from more than one provider (fusion) or ties a
provider's coordinate to an asset defined in another provider's frame (hand-off, or a
surveyed landmark), the \emph{interop tax} is the position error induced by
$\{\datum_{ab},\resid_{ab}\}$ that a single self-consistent provider never sees.
\end{definition}

\begin{definition}[Reducible and irreducible components]
\label{def:redirr}
Collect the pairwise disagreements \eqref{eq:diffmodel}. The recovered Moon rotation
admits the measured affine split
$\rot_{\mathrm{moon}}=\rot_{\mathrm{tie}}+\rot_{\mathrm{exc}}$, and the cross-provider
error separates into three terms:
\begin{itemize}
  \item a \emph{reducible} part $\reduc$: the Helmert rotation
  $\rot_{\mathrm{tie}}$ that is common across bodies and directions and is therefore
  removable by adopting a common frame realization---a coordinate convention.
  Operationally we estimate it as the body-common (planet-common) frame-tie rotation
  shared by the Moon and the planetary positions;
  \item an \emph{irreducible} part $\irr$: the body-specific Moon-\emph{excess}
  rotation $\rot_{\mathrm{exc}}=\rot_{\mathrm{moon}}-\rot_{\mathrm{tie}}$, a
  lunar-orbit orientation difference that does \emph{not} appear in the planetary
  frame-tie. No coordinate convention removes $\irr$; only a common designated
  dynamical reference (a common ephemeris) sets $\irr\to0$ by construction;
  \item a non-Helmert \emph{residual} $\resid_{ab}$: the part of the difference that
  no constant rigid transform absorbs. It is bookkept separately from $\rot_{\mathrm{exc}}$
  (it is small, sub-nanoradian in rotation-equivalent), is likewise irreducible by any
  coordinate convention, and enters the budget total as its own term
  (Table~\ref{tab:budget}).
\end{itemize}
The split $\reduc$/$\irr$ is affine (not orthogonal): $\rot_{\mathrm{tie}}$ and
$\rot_{\mathrm{exc}}$ are in general non-parallel, so their metre envelopes do not add
in quadrature to the raw disagreement (Section~\ref{sec:budget}).
\end{definition}

\noindent The reducible/irreducible split is the load-bearing structure of the
paper. Section~\ref{sec:theorem} proves that a common frame convention removes only
the body-common $\reduc$ and leaves $\irr$ (and the residual) intact;
Section~\ref{sec:interephemeris} measures both parts on the real ephemerides; and
Section~\ref{sec:budget} inverts the budget to the tolerance $\tau(B)$. We flag now,
and repeat where it matters, that the \emph{attribution} of the planet-common
rotation to a reducible frame-tie (as opposed to shared dynamics) is a stated
modelling interpretation; the convention-free claim is the Moon-excess rotation
(Definition~\ref{def:redirr}, second bullet), which is unambiguously Moon-specific and
hence not removable by any whole-frame choice.
